% !PRINTED 660
\weberpage{660}

Der Satz wird bewiesen sein, wenn von irgend einer Potenz
$\varphi^k$ von $\varphi$, deren Exponent nicht durch $q$ theilbar
ist, bewiesen werden kann, dass sie der Hauptclasse angeh\"ort.
Denn sind $\alpha, \beta$ Zahlen in $\Omega_m$ und ist
\[
\varphi^m = \varepsilon' \alpha, \quad \varphi^k = \varepsilon'' \beta,
\]
so k\"onnen wir die ganzen rationalen Zahlen $x, y$ so bestimmen,
dass
\[
mx + ky = 1
\]
wird, und dann wird
\[
\varphi = \varphi^{mx} \varphi^{ky} = \varepsilon''' \alpha^x \beta^y,
\]
also ist auch, wenn $\varepsilon', \varepsilon'', \varepsilon'''$
Einheiten sind, $\varphi$ selbst mit einer Zahl associirt.

Nach dem heutigen Stande der Sache ist der Beweis f\"ur ein
gerades $m$ auf einem ganz anderen Wege zu f\"uhren, als f\"ur
ein ungerades, und wir besch\"aftigen uns also zun\"achst mit dem
leichteren Falle des ungeraden $m$ $^1$).

Die in a) \"uber $\varphi_n$ gemachten Voraussetzungen k\"onnen
wir, wenn $\alpha, \beta$ wieder irgend welche Zahlen des
K\"orpers $\Omega_m$ und $\varepsilon$ Einheitsfunctionale
bedeuten, durch die beiden Formeln ausdr\"ucken:
\weq{1}{\varphi_n^m = \varepsilon \beta,}
\weq{2}{\varphi_n = \varepsilon \alpha \varphi_1^n.}
Wir werden nun die Aufgabe schrittweise l\"osen.

Ist zun\"achst, wie fr\"uher, $\sigma$ der in $q$ enthaltene
Primfactor, der eine in $\Omega_m$ existirende Zahl ist, und
$\sigma^k$ die in $\varphi$ enthaltene Potenz von $\sigma$, so
setzen wir
\weq{3}{\varphi = \sigma^k \psi,}
worin $\psi$ ein zu $q$ theilerfremdes Functional bedeutet, das
denselben Bedingungen (1), (2) wie $\varphi$ gen\"ugt, und wenn
eines von beiden in die Hauptclasse geh\"ort, so gilt das auch
vom anderen. Der Satz a) braucht also nur noch unter der
Voraussetzung bewiesen zu werden, dass $\varphi$ theilerfremd zu
$q$ ist.

Es sei nun $p$ eine von $q$ verschiedene Primzahl, die in
Primfactoren zerlegt den Ausdruck hat:
\weq{4}{p = \prod^{\xi} \mathfrak{p}_\xi,}
%CONT>>

\noindent
$^1$) Vgl. \"ubrigens den neuen Beweis von \gesperrt{Hilbert}, bei
dem der Unterschied zwischen geradem und ungeradem $m$ weit mehr
zur\"ucktritt (S. 648).
